\section{A Common Representation: Position Automata}

We use position automata as the bridge representation.  A regular expression is
first transformed into a positioned expression in which every atom occurrence is
assigned a fresh natural-number position.  For example, \(ab^{*}\) becomes
\(a_1 b_2^{*}\).  The construction computes the usual nullable, firstpos,
lastpos, and followpos sets.

The Coq development defines source expressions and positioned expressions as:
\[
\begin{array}{rcl}
r &::=& \emptyset \mid \epsilon \mid a \mid r+r \mid rr \mid r^{*}\\
\hat r &::=& \emptyset \mid \epsilon \mid (p,a) \mid \hat r+\hat r
        \mid \hat r\hat r \mid \hat r^{*}.
\end{array}
\]
The labeling function \(\mathsf{label}\) traverses the source syntax
left-to-right and assigns consecutive positions.  The automaton states are sets
of positions.  The start state is \(\mathsf{firstpos}(\hat r)\), the transition
from a state \(S\) on symbol \(a\) is the union of
\(\mathsf{followpos}(p)\) for positions \(p\in S\) whose label matches \(a\),
and finality is determined by intersection with \(\mathsf{lastpos}(\hat r)\),
with a nullable special case for the empty word.

The development also defines a marked-word semantics.  A marked word records not
only the consumed symbols, but the positions that consumed them:
\[
  [(p_1,a_1),\ldots,(p_n,a_n)].
\]
Marked words are a useful proof device because they expose the exact syntactic
occurrences used by a match.  The theorem already mechanized in the development
shows that labeling preserves Kleene semantics after erasing positions.
